\section{Discusión y Conclusiones}\label{sec:conclusiones}
Los cinco modelos terminan en el mismo lugar: AUC de 0,75 a 0,76 en validación y de 0,77 en prueba, con intervalos que se traslapan. Que una regresión logística iguale a una SVM con kernel o a una red neuronal indica que el límite lo pone la información disponible y no la capacidad del modelo. Once preguntas de encuesta, con el IMC agrupado en cuatro categorías, no alcanzan a describir el riesgo con más detalle. El resultado está dentro del rango de Xie et al. \cite{xie2019} con la misma encuesta; los AUC cercanos a 0,99 de otros trabajos \cite{majcherek2025} usan más variables, entre ellas la salud general percibida, y otra preparación de los datos, así que no son directamente comparables.

Donde sí hay diferencias es en la métrica clave. Con la ponderación de clases, todos los modelos detectan entre 74\% y 80\% de los casos en prueba, frente al 51,6\% de Xie et al.; la SVM es la que más detecta, a cambio de más falsas alarmas. SMOTE no mejoró frente a la ponderación de clases y SMOTEN empeoró los resultados, lo que confirma que, con un desbalance moderado como este, ponderar las clases basta. También quedó claro que la codificación importa más de lo que parece: tratar las ordinales como enteros ayuda a los modelos de distancia pero perjudica a la regresión logística, porque el efecto de la edad no es proporcional, y conservar ``no reportado'' como categoría tiene sentido porque no responder se asocia con la diabetes en seis de las diez variables con faltantes.

La reducción de dimensión deja una conclusión práctica y otra metodológica. La práctica es que sexo, frutas y vegetales se pueden retirar sin perder desempeño: un cuestionario de tamizaje con ocho preguntas en lugar de once es más corto de aplicar, y eso tiene más valor en este contexto que ahorrar columnas. La metodológica es que las técnicas de extracción rinden muy distinto según la codificación. PCA no reduce de verdad la regresión logística, porque sus 45 componentes son casi el rango completo de la codificación \textit{one-hot}; en la SVM, que recibe las ordinales como números, 20 componentes conservan todo el desempeño y 10 pierden solo un punto de AUC, aunque a costa de que las nuevas variables ya no se pueden interpretar. UMAP no funcionó con estos datos categóricos, y su versión supervisada mostró que un paso de reducción también puede sobreajustar si se le da la etiqueta, algo que solo se detecta porque se evaluó dentro de la validación cruzada.

Con todo esto, la regresión logística sin sexo, frutas y vegetales es la opción más recomendable para un tamizaje: tiene el mismo AUC que cualquier otro modelo (0,77), detecta tres de cada cuatro casos, se entrena en un segundo con todos los datos y sus coeficientes se pueden explicar. Si la prioridad es detectar más casos aunque aumenten las falsas alarmas, la SVM llega a una sensibilidad de 0,80.

Este trabajo tiene limitaciones que conviene tener presentes. El diagnóstico es autorreportado, así que entre los negativos hay personas con diabetes no diagnosticada; los datos son transversales, de modo que el modelo identifica perfiles asociados con tener diabetes y no predice quién la desarrollará; no se usaron los pesos muestrales de la encuesta, y faltan cuatro jurisdicciones. Como trabajo futuro quedan ajustar el umbral de decisión para fijar una sensibilidad objetivo, calibrar las probabilidades e incorporar variables como la salud general percibida, la hipertensión o el colesterol, que el BRFSS también recoge.

